\scanpage{60}
% manuscript line(s) 1
\[
\underbrace{0\subset L^{\otimes2}}_{L^{\otimes2}}\ \underbrace{\subset L\otimes E}_{L^{\deleted{\unsure{\otimes2}}}}\ \underbrace{\subset\mathrm{Sym}^2(E)}_{\underline O_X}
\]

% manuscript line(s) 2
Considérons \deleted{\unsure{écrivons}}
% manuscript line(s) 3
\[
Q(x)-x^2:\quad\unsure{P|E}\mathbin{-}\mathrm{Sym}^2(E)
\]
% manuscript line(s) 4–6
fonction \deleted{quadratique, sa partie homogène est} \inserted{polynomiale de degré \deleted{\unsure{possiblement}} $\leq2$, dont la} partie homogène \deleted{\unsure{est sa}} quadratique
% manuscript line(s) 6
\[
\deleted{Q(x+u)=}\qquad L\mathbin{-}L^{\otimes2}
\]
% manuscript line(s) 7
est nulle. Donc c'est une application \uline{affine}

% manuscript line(s) 8–12
\deleted{soit $Q(x)=x^2+\mathrm{Af}(x)$ \quad $\mathrm{Af}:P\mathbin{-}\deleted{\unsure{E}}\mathrm{Sym}^2 E$ application \uline{affine}. Son terme linéaire est un hom $L\longrightarrow\deleted{\unsure{E}}\mathrm{Sym}^2(E)$, \inserted{donc} défini par un \unsure{tel} élément $T$ de $L^\vee\otimes\mathrm{Sym}^2 E\supset E$}

% manuscript line(s) 13–15
Utilisant un scindage de \unsure{$N$ cer. $E$|l'ext. $E$|$N$ et $E$} par une section $e$ de $P_1$, i.e. $e\in\Gamma E$, $\psi(e)=1$, on aura, posant
% manuscript line(s) 15
\[
\begin{cases}x=e+u\\ u=x-e\end{cases},\quad\inserted{u\in\Gamma\underline L,}\quad\text{et écrivant}
\]
% manuscript line(s) 16
\[
Q(e+u)=u^2+bu+c\qquad\begin{aligned}u&\in\Gamma L\\ b&\in\Gamma\unsure{L'|L}\\ c&\in\Gamma L^{\otimes2}\end{aligned}
\]
% manuscript line(s) 17–18
\[
\begin{aligned}Q(\underbrace{e+u}_{x})-\deleted{\unsure{Q}}(\underbrace{e+u}_{x})^2&=bu+c-2eu-e^2=\underbrace{u}_{x-e}(b-2e)+c-e^2\\ &=\underline x(\unsure{+}b-2e)\deleted{\unsure{-e(b-e)+c-e^2}}+(e^2+c-be)\end{aligned}
\]
% manuscript line(s) 19–21
\[
\text{i.e.}\quad\left\{\begin{aligned}Q(x)&=x^2\deleted{\unsure{+}}\unsure{-|+}xT+N\\[-2pt] &\text{où}\quad\boxed{T=2e-b\in E,\quad N=e^2+\underset{\substack{\in\\\Gamma L^2}}{c}-\underset{\substack{\in\\\Gamma(L\otimes E)\subset\Gamma(\mathrm{Sym}^2E)}}{be}\in\mathrm{Sym}^2(E)}\\ &\text{avec}\quad\boxed{\psi(T)=2,\varepsilon(N)=1}\end{aligned}\right.
\]
% manuscript line(s) 21
\diagram{\begin{tikzcd}[ampersand replacement=\&,column sep=small]0\arrow[r,dash]\&L\otimes E\arrow[r]\&\mathrm{Sym}^2(E)\arrow[r,dash,"\varepsilon"]\&\underline O_X\arrow[r,dash]\&0\end{tikzcd}}

% manuscript line(s) 22–26
\deleted{\illegible{2}} On vérifie que $T$, donc aussi $N$, ne dépend pas du scindage choisi. [$T$ est élém\textsuperscript{t} d'un torseur \inserted{$P_2=\psi^{-1}(2)$} sous $L^{\deleted{\unsure{\otimes}}}$, \deleted{\unsure{donc|et}} et $N$ (pour $T$ choisi), \deleted{que} la condition que $N$ soit à valeurs dans $L^{\otimes2}$, est élém\textsuperscript{t} d'un torseur sous $L^{\otimes2}\subset\mathrm{Sym}^2(E)$, donc \inserted{le choix de} $Q\deleted{\unsure{+2}}\longleftrightarrow(T,N)$ est \unsure{celui}
% archive sheet number: 59
